\documentclass{article}
\usepackage{tcee-repaso}
\tcetema[El SEI desde Bretton Woods]{3.B.21}{El sistema económico internacional desde la desaparición del sistema de Bretton-Woods.}
\tcefecha{13/11/2026}
\begin{document}
\cabecerarepaso
\begin{columnas}

% ==========================================================================
\bloque{Esquema del cante}
\esquemacante{%
\apartadocante{Introducción: «lo provisional es lo que más dura»}{2,5}
\apartadocante{I. Del fin de Bretton Woods a la Gran Moderación}{11}
\subcante{1. La ruptura (1971-1976)}{2}
\subcante{2. Los primeros años del no sistema (1973-1979)}{1,5}
\subcante{3. Gran Moderación: sistema comercial}{2,5}
\subcante{4. Gran Moderación: monetario y financiero}{5}
\apartadocante{II. De la crisis financiera global a la pandemia}{7}
\subcante{1. Crisis financiera global y crisis del euro}{2}
\subcante{2. Consecuencias sobre el sistema}{3,5}
\subcante{3. La COVID-19}{1,5}
\apartadocante{III. La actualidad: retos y reformas}{8,5}
\subcante{1. Comercio: fragmentación y aranceles}{3}
\subcante{2. Monetario y financiero: el dólar y sus rivales}{3}
\subcante{3. Retos y propuestas de reforma}{2,5}
\apartadocante{Conclusión}{1}}
\emph{Gráficos del cante}: el trilema con el paso de Bretton Woods a la flotación (\hyperref[rep:I]{I}) y el circuito de Bretton Woods II (\hyperref[rep:I.4]{I.4}); función clave: $\mathrm{CC}\equiv S-I$. En cada etapa, los 3 planos en el mismo orden: comercial, monetario, financiero.

% ==========================================================================
\bloque[2,5]{Introducción}
\begin{itemize}
\item \textbf{Enganche}: «Lo provisional es lo que más dura»: \autor{Nixon} suspende «temporalmente» la convertibilidad (15/8/1971) y, 55 años después, ni el dólar ha vuelto al oro ni hay otro sistema. Abril de 2025: aranceles «recíprocos»; arancel efectivo de EE.\,UU. 2,6\,\% (2024) $\to$ 8\,\% (2025), máximo desde 1946. Y el dólar, 56,7\,\% de las reservas y 89\,\% de las operaciones de cambio. \emph{¿Hay un sistema o solo un «no sistema» que sobrevive porque nadie ha construido otro?}
\item \textbf{Relevancia}: las reglas deciden a qué tipo se comercia, quién se ajusta y quién presta en las crisis; la historia, pilar del análisis (\autor{Schumpeter}; «esta vez es distinto», \autor{Reinhart}-\autor{Rogoff}); España: SME 1992-1993, euro, política comercial de la UE \vertema{4.A.21}.
\item \textbf{Contextualización}: \concepto{SEI}: instituciones, normas y prácticas entre países, en 3 \emph{planos}: comercial (GATT/OMC) \vertema{3.B.33}, monetario (régimen, reservas, liquidez) \vertema{3.B.15}, financiero (capitales, regulación, crisis). 3 funciones del SMI: \emph{ajuste}, \emph{liquidez}, \emph{confianza}. Bretton Woods (1944): 3 \emph{pilares} (FMI, BIRF, GATT), cambios oro a 35\,\$/onza, controles; nada dura para siempre: \autor{Triffin} \vertema{3.B.20}. \concepto{No sistema} (\autor{Williamson}, 1976): sin reglas comunes ni activo de reserva definido.
\item \textbf{Problemática}: ¿por qué cayó y qué lo sustituyó? ¿Cómo se han cumplido ajuste, liquidez y confianza? ¿Qué cambiaron las crisis? ¿Dónde estamos y qué reformas?
\item \textbf{Estructura}: orden cronológico, y en cada etapa los 3 planos: \hyperref[rep:I]{I} 1971-2007; \hyperref[rep:II]{II} 2007-2021; \hyperref[rep:III]{III} actualidad.
\end{itemize}

% ==========================================================================
\bloque[11]{I. Del fin de Bretton Woods a la Gran Moderación\label{rep:I}}
\begin{ideaclave}
Cae porque el capital se mueve y el dólar no puede ser reserva y convertible a la vez. Lo sustituye un no sistema de libertad cambiaria: hiperglobalización, estabilidad en los avanzados, crisis en los emergentes con tipo fijo y, al final, Bretton Woods II.
\end{ideaclave}

\sub{I.1. La ruptura (1971-1976)\label{rep:I.1}}
\begin{itemize}
\item \emph{Por qué}: ajuste asimétrico (solo deficitarios); liquidez con déficits de EE.\,UU. que minan la convertibilidad (\autor{Triffin}, 1960); confianza (pasivos $>$ oro desde 1960; Vietnam). \emph{Trilema} (\autor{Fleming}, 1962; \autor{Mundell}, 1963): con capital libre, de tipo fijo + controles a flotación \vertema{3.B.15}.
\item \emph{4 pasos}: \textbf{ago. 1971} fin de la convertibilidad + recargo del 10\,\%; \textbf{dic. 1971} Smithsonian: 38\,\$/onza, bandas $\pm1\to\pm2{,}25$\,\%; serpiente (1972); \textbf{mar. 1973} flotación; \textbf{ene. 1976} Jamaica (Kingston), Segunda Enmienda en vigor el 1/4/1978: libertad cambiaria (salvo oro), vigilancia del FMI (art. IV), desmonetización del oro, DEG como aspiración.
\end{itemize}
\graficorepaso[0.72\linewidth]{trilema-ruptura}{Trilema: de Bretton Woods (tipo fijo con controles) a la flotación, 1971-1973}

\sub{I.2. Primeros años del no sistema (1973-1979)\label{rep:I.2}}
\begin{itemize}
\item \textbf{Comercial}: Ronda de Tokio (1973-1979): aranceles $-1/3$, fórmula general (suiza, $t_1=at_0/(a+t_0)$), «códigos» no arancelarios; nuevo proteccionismo (restricciones voluntarias, Multifibras 1974).
\item \textbf{Monetario}: flotación sucia sin cooperación; dólar ↓ hasta 1978; serpiente y SME (1979) \vertema{3.B.16}; Bonn 1978, locomotora \vertema{3.B.19}.
\item \textbf{Financiero}: reciclaje de petrodólares hacia América Latina; Herstatt (1974) $\to$ Comité de Basilea (fin de 1974, G-10) \vertema{3.B.27}.
\end{itemize}

\sub{I.3. Gran Moderación: sistema comercial\label{rep:I.3}}
\begin{itemize}
\item \concepto{Hiperglobalización}: comercio 37\,\% del PIB mundial (1985) $\to$ 60\,\% (2008).
\item Ronda Uruguay (1986-1994): agricultura, textiles, servicios, propiedad intelectual; Marrakech: \concepto{OMC} (1/1/1995), solución de diferencias con apelación \vertema{3.B.33}; Doha (2001).
\item Tigres (exportación) vs. América Latina (sustitución); China (apertura 1978; OMC 11/12/2001; choque de China, \autor{Autor}-\autor{Dorn}-\autor{Hanson}, 2013); muro de Berlín (1989).
\item \concepto{CGV}: contenedor y TIC, «segunda separación» (\autor{Baldwin}) \vertema{3.B.18}; regionalismo (mercado único, TLCAN, Mercosur, ASEAN) \vertema{3.B.35}.
\item Curva del elefante (\autor{Lakner}-\autor{Milanovic}, 2016): ganan clases medias asiáticas y el 1\,\% superior.
\end{itemize}
\graficorepaso[0.9\linewidth]{apertura-comercial}{Comercio mundial, \% del PIB: hiperglobalización hasta 2008 y lentiglobalización después}

\sub{I.4. Gran Moderación: monetario y financiero\label{rep:I.4}}
\begin{itemize}
\item \textbf{Tipo flexible, cooperación}: \autor{Volcker} (1979-1982) + fiscal expansiva: dólar real +50\,\% hasta mar. 1985; \textbf{Plaza} (22/9/1985, G-5): depreciar; \textbf{Louvre} (22/2/1987, G-6, dólar ya $-23$\,\%): zonas objetivo secretas \vertema{3.B.19}. Desde los 90, ancla = objetivos de inflación (Nueva Zelanda, 1990). Euro (1999; 21 países desde 2026).
\item \textbf{Tipo fijo, crisis} \vertema{3.B.17}: \emph{1982} (petrodólares, \autor{Volcker}; México, ago. 1982; década perdida; Baker 1985, Brady 1989) \vertema{3.B.29}; \emph{SME 1992-1993} (2.ª generación; bandas $\pm15$\,\%); \emph{México 1994-1995} (mixta, \emph{tesobonos}, parada súbita; PIB $-5{,}9$\,\%); \emph{Asia 1997-1998} (gemelas, riesgo moral; Indonesia $-13{,}1$\,\%); \emph{Rusia 1998}, \emph{Brasil 1999}; \emph{Argentina 2001-2002} (caja de conversión, corralito; PIB $-10{,}9$\,\%).
\item Respuesta: Foro de Estabilidad Financiera y G-20 de ministros (1999); Basilea I (1988) y Basilea II (2004); reservas como seguro.
\end{itemize}
\graficorepaso[0.95\linewidth]{dolar-tipo-efectivo}{Dólar real (BPI): máximo de mar. 1985, Plaza y Louvre ($-23$\,\%), mínimo de 2011 y máximo de 2022}
\textbf{Bretton Woods II} (\autor{Dooley}, \autor{Folkerts-Landau} y \autor{Garber}, 2003):
\begin{itemize}
\item \emph{a) Idea}: centro (EE.\,UU., déficit, activo seguro) y periferia (Asia, China) con moneda infravalorada: Bretton Woods sin tratado.
\item \emph{b) i) Supuestos}: periferia con trabajo excedente y finanzas ineficientes que crece exportando; interviene y controla capitales; el centro emite el activo seguro; identidades contables. \emph{ii) Desarrollo}:
\begin{formula}
\centering \funcionclave{\mathrm{CC}\equiv S-I}
\end{formula}
\item[] 1) la periferia exporta ($S>I$); 2) compra dólares, $\Delta R>0$; 3) compra deuda del Tesoro: financia el déficit del centro ($S<I$) a tipos bajos; 4) inversión directa del centro en la periferia, con las reservas como garantía. \emph{iii) Implicaciones}: ganan los 2; estable mientras dure el excedente de trabajo.
\item \emph{c) Extensiones}: exceso de ahorro (\autor{Bernanke}, 2005); escasez de activos seguros; privilegio exorbitante (\autor{Gourinchas}-\autor{Rey}, 2007); \autor{Triffin} moderno (\autor{Farhi}-\autor{Gourinchas}-\autor{Rey}, 2011).
\item \emph{d) Evidencia}: déficit de EE.\,UU. $-1{,}5$\,\% del PIB mundial (2006); China $+0{,}65$\,\% (2008) y $+0{,}67$\,\% (2026p); reservas 1,8 (1999) $\to$ $>$6 billones (2007).
\item \emph{e) Valoración}: explicó la duración del déficit; pero la periferia no es un bloque, la estabilidad alimentó la burbuja (\autor{Obstfeld}-\autor{Rogoff}, 2009) y China flexibilizó en 2005.
\end{itemize}
\graficorepaso[0.85\linewidth]{bretton-woods-ii}{Bretton Woods II: exportaciones, reservas, deuda del Tesoro e inversión directa}
\graficorepaso[0.92\linewidth]{desequilibrios-globales}{Desequilibrios globales (\% del PIB mundial): EE.\,UU. $-1{,}5$ (2006); China $+0{,}67$ (2026p)}

% ==========================================================================
\bloque[7]{II. De la crisis financiera global a la pandemia\label{rep:II}}
\begin{ideaclave}
2 crisis mundiales muestran los límites del no sistema y su capacidad de reacción: G-20, nueva regulación y la Fed como prestamista del mundo; el comercio deja de tirar y se politiza.
\end{ideaclave}

\sub{II.1. Crisis financiera global y del euro\label{rep:II.1}}
\begin{itemize}
\item \emph{2007-2009} \vertema{3.B.26}: \emph{subprime} titulizado, supervisión laxa, banca en la sombra, ahorro exterior; Fed 1\,\% (2004) $\to$ 5,25\,\% (jun. 2006); interbancario (ago. 2007); Lehman (15/9/2008); PIB mundial $-0{,}4$\,\% (2009).
\item \emph{Euro 2010-2012}: bancos-soberanos, sin activo seguro ni prestamista; MEDE, PEC reforzado, desequilibrios, Unión Bancaria, «lo que sea necesario» (jul. 2012); Grecia, corralito (2015) \vertema{3.B.16}.
\end{itemize}

\sub{II.2. Consecuencias\label{rep:II.2}}
\begin{itemize}
\item \textbf{Comercial}: gran colapso (2009, $>$10\,\% en volumen), sin proteccionismo general; \emph{lentiglobalización}; Doha atascada (Facilitación, 2013); Órgano de Apelación bloqueado (dic. 2019), MPIA (62 miembros, 60\,\% del comercio); acuerdos profundos (UE-Canadá, CPTPP, RCEP, T-MEC); guerra EE.\,UU.-China (2018-2020; fase uno, ene. 2020).
\item \textbf{Financiero}: G-20 de líderes, foro principal (Pittsburgh 2009): $\approx$85\,\% PIB, 75\,\% comercio, 2/3 población; \textbf{Londres (2/4/2009)}: FMI ×3 hasta 750.000 M\$, DEG 250.000 M\$, FSF $\to$ \textbf{FSB}; demasiado grandes para caer, resolución, derivados, no bancaria; \textbf{Basilea III} (2010-2017) \vertema{3.B.27}.
\item \textbf{Monetario}: \emph{swaps} de la Fed (permanentes con 5 bancos centrales, 2013); ciclo financiero global, «más que un trilema, un dilema» (\autor{Rey}, 2013) \vertema{3.B.15}; \emph{taper tantrum} (2013); guerras de divisas, G-7 (2013).
\end{itemize}

\sub{II.3. La COVID-19\label{rep:II.3}}
\begin{itemize}
\item PIB mundial $-2{,}7$\,\% (2020); salida récord de capital de los emergentes.
\item Respuesta: fiscal (CARES, \emph{Next Generation EU}); Fed: \emph{swaps} a 9 más y FIMA (permanente, jul. 2021); FMI: DEG 456.000 M (650.000 M\$, ago. 2021) \vertema{3.B.22}; ISSD ($>$12.000 M\$, $\approx$50 países) y Marco Común (nov. 2020) \vertema{3.B.29}.
\item Lección: CGV frágiles $\to$ resiliencia, autonomía estratégica, política industrial.
\end{itemize}

% ==========================================================================
\bloque[8,5]{III. La actualidad: retos y reformas\label{rep:III}}
\begin{ideaclave}
Desde 2022 el sistema se fragmenta: seguridad antes que eficiencia, el antiguo garante del multilateralismo lo cuestiona y los bloques se ven en comercio, reservas y pagos. Aun así, el dólar sigue en el centro y el comercio resiste.
\end{ideaclave}

\sub{III.1. Comercio: fragmentación y aranceles\label{rep:III.1}}
\begin{itemize}
\item \emph{Fragmentación}: Ucrania (2022), sanciones, $\approx$300.000 M\$ de reservas rusas congeladas; IRA y CHIPS (2022); \emph{friend-shoring}; coste hasta 7\,\% del PIB mundial (FMI, 2023).
\item \emph{Aranceles de EE.\,UU.}: IEEPA (fentanilo, feb.-mar. 2025; «recíprocos» 2/4/2025, 10\,\% universal); China hasta 145\,\%, tregua de Ginebra (may. 2025) y Busan (oct.-nov. 2025), hasta ene. 2027; acuerdos: Reino Unido 10\,\%, UE-Japón-Corea techo 15\,\%; sección 232: acero y aluminio 50\,\%, automóviles 25\,\%. \textbf{Supremo, 20/2/2026} (\emph{Learning Resources}), 6-3: la IEEPA no permite aranceles. Sustitución: sección 122 (10\,\%, feb.-jul. 2026) y sección 301 por «trabajo forzoso» (10-12,5\,\%, 60 economías, desde 24/7/2026); devoluciones.
\item \emph{OMC}: NMF 80 $\to$ 72\,\% del comercio; CM14 (Yaundé, mar. 2026) sin acuerdo: expira la moratoria del comercio electrónico; 166 miembros.
\item \emph{Resistencia}: comercio de mercancías +4,6\,\% (2025), +3,9\,\% (2026p, OMC); desvío vía Vietnam y México. \emph{Choque}: guerra con Irán (desde 28/2/2026) y cierre de Ormuz; PIB mundial 3,1\,\% (2026p).
\end{itemize}
\graficorepaso[0.9\linewidth]{arancel-eeuu}{Arancel efectivo de EE.\,UU.: 18,7\,\% (1933), 1,6\,\% (2017), 8\,\% (2025)}

\sub{III.2. Monetario y financiero: el dólar y sus rivales\label{rep:III.2}}
\begin{itemize}
\item \emph{Dominancia del dólar}: reservas 71\,\% (1999) $\to$ 56,7\,\% (jun. 2026); euro 20,6\,\%; renminbi 2,1\,\%; mercado de divisas 89,2\,\% (2025); facturación del comercio: moneda dominante (\autor{Gopinath}, 2020) \vertema{3.B.13}.
\item \emph{Diversificación}: oro $>$1.000 t/año desde 2022; 27\,\% de las reservas a precios de mercado (2025) vs. Tesoro 22\,\% y euro 15\,\% (BCE); monedas no tradicionales.
\item \emph{Renminbi}: desde 2009; DEG (2016, 12,28\,\%), \emph{swaps}, CIPS, mBridge; sin convertibilidad de la cuenta de capital no hay reserva mundial.
\item \emph{Dinero nuevo}: ley GENIUS (jul. 2025, \emph{stablecoins} en dólares); euro digital (mandato del Parlamento, jul. 2026; piloto 2027).
\item \emph{Desequilibrios}: vuelven a crecer; el FMI pide reequilibrar con ahorro e inversión, no con aranceles ($\mathrm{CC}\equiv S-I$).
\item \emph{Red de seguridad}: reservas, \emph{swaps}, regionales (Chiang Mai, MEDE) y FMI: amplia pero desigual.
\end{itemize}
\graficorepaso[0.9\linewidth]{reservas-monedas}{Reservas de divisas por monedas: el dólar baja de 71 a 56,7\,\%; el renminbi, 2,1\,\%}

\sub{III.3. Retos y propuestas de reforma\label{rep:III.3}}
\begin{itemize}
\item \emph{Gobernanza}: FMI, 16.ª revisión (+50\,\%, dic. 2023) sin entrar en vigor (Congreso de EE.\,UU.; plazo 15/11/2026); 17.ª sin acuerdo \vertema{3.B.22}; G-20 dividido (Johannesburgo 2025 sin EE.\,UU.; Miami, dic. 2026); EE.\,UU. sale de la OMS, de París (27/1/2026) y de 66 organizaciones, no del FMI.
\item \emph{Monetario}: red de seguridad (recursos, \emph{swaps}, reciclar DEG); multipolaridad (euro, deuda común); DEG supranacional (\emph{bancor}, \autor{Triffin}).
\item \emph{Comercial}: reforma de la OMC, plurilaterales, acuerdos profundos (UE-Mercosur, dic. 2024) \vertema{3.B.41}.
\item \emph{Financiero}: Basilea III (UE desde 2025, cartera de negociación 2027; EE.\,UU., propuesta mar. 2026), no bancaria, cripto.
\item \emph{Fiscal}: pilar 2, mínimo del 15\,\% ($\geq$750 M€); G-7 (jun. 2025) y paquete «lado a lado» (ene. 2026): exime a los grupos de EE.\,UU. \vertema{4.B.18}.
\item \emph{Clima}: París sin EE.\,UU.; CBAM; clubes climáticos (\autor{Nordhaus}) \vertema{3.B.31}.
\end{itemize}

% ==========================================================================
\bloque[1]{Conclusión}
\begin{itemize}
\item \textbf{Recapitulación}: «En conclusión,» ruptura 1971-1973 y no sistema legalizado en Jamaica; Gran Moderación: hiperglobalización, crisis emergentes, Bretton Woods II; 2008 y 2020 superadas con G-20, regulación y Fed; desde 2022, fragmentación con el dólar en el centro.
\item \textbf{Relevancia}: España, abierta y en el euro: mercados, financiación y crisis.
\item \textbf{Extensiones}: \vertema{3.B.20} \vertema{3.B.22} \vertema{3.B.33} \vertema{3.B.19} \vertema{3.B.17} \vertema{3.B.26} \vertema{3.B.15} \vertema{3.B.29} \vertema{3.B.32}.
\item \textbf{Opinión}: resistente no por sus reglas, sino por la cooperación de los bancos centrales y la confianza en el dólar; si falla uno, no hay red de repuesto.
\item \textbf{Idea final}: «lo provisional es lo que más dura»: la suspensión «temporal» cumplió 55 años; hay un orden de hecho, y la pregunta es si sobrevivirá a la fragmentación o habrá que construir uno nuevo sin diseñarlo.
\end{itemize}

% ==========================================================================
\bloque{Autores y fechas}
\hito{1953}{\autor{Friedman}: a favor de los tipos flexibles}
\hito{1960}{\autor{Triffin}: el dilema}
\hito{1962}{\autor{Fleming}; 1963, \autor{Mundell}: trilema}
\hito{1971}{Fin de la convertibilidad; Smithsonian}
\hito{1973}{Flotación; Ronda de Tokio}
\hito{1974}{Comité de Basilea}
\hito{1976}{Jamaica; \autor{Williamson}: «no sistema»}
\hito{1979}{SME; \autor{Volcker}}
\hito{1982}{Crisis de la deuda}
\hito{1985}{Plaza; 1987, Louvre}
\hito{1989}{Plan \autor{Brady}}
\hito{1995}{OMC}
\hito{1999}{Euro; G-20 de ministros; \autor{Krugman}: tercera generación}
\hito{2001}{China en la OMC; Doha}
\hito{2003}{\autor{Dooley}, \autor{Folkerts-Landau} y \autor{Garber}: Bretton Woods II}
\hito{2005}{\autor{Bernanke}: exceso de ahorro}
\hito{2009}{Londres; FSB}
\hito{2013}{\autor{Rey}: dilema}
\hito{2016}{\autor{Lakner}-\autor{Milanovic}: elefante; renminbi en el DEG}
\hito{2020}{\autor{Gopinath}: moneda dominante}
\hito{2023}{FMI: fragmentación, hasta 7\,\%}
\hito{2025}{Aranceles «recíprocos»; ley GENIUS}
\hito{2026}{Sentencia IEEPA; CM14; sección 301; Ormuz}

% ==========================================================================
\bloque{Datos clave}
\begin{itemize}
\item Arancel efectivo de EE.\,UU.: \dato{2,6\,\% (2024) $\to$ 8\,\% (2025)}{BEA, 2026}.
\item Reservas: dólar \dato{56,7\,\%}{FMI, COFER, jun. 2026}; euro \dato{20,6\,\%}{FMI}; renminbi \dato{2,1\,\%}{FMI}.
\item Mercado de divisas: dólar \dato{89,2\,\%}{BPI, 2025}; renminbi \dato{8,5\,\%, 5.ª}{BPI}.
\item Oro: \dato{27\,\% de las reservas oficiales (2025)}{BCE, 2026}.
\item Comercio mundial: \dato{37\,\% (1985) $\to$ 60\,\% del PIB (2008)}{BM}; volumen \dato{+4,6\,\% (2025), +3,9\,\% (2026p)}{OMC}.
\item Desequilibrios: EE.\,UU. \dato{$-1{,}5$\,\% PIB mundial (2006)}{FMI}; China \dato{$+0{,}67$\,\% (2026p)}{FMI}.
\item PIB mundial: \dato{$-0{,}4$\,\% (2009), $-2{,}7$\,\% (2020), 3,1\,\% (2026p)}{FMI, WEO oct. 2026}.
\item Dólar real: \dato{$-23$\,\% de mar. 1985 a feb. 1987}{BPI}.
\item OMC: \dato{166 miembros}{OMC}; MPIA \dato{62, 60\,\% del comercio}{CE, 2026}.
\end{itemize}

% ==========================================================================
\bloque{Preguntas del tribunal}
\begin{enumerate}
\item \textbf{¿Por qué cayó Bretton Woods?} Ajuste asimétrico, \autor{Triffin} y confianza; con el capital libre, el trilema obliga a soltar el tipo fijo.
\item \textbf{¿Qué decidió Jamaica?} Libertad cambiaria, vigilancia del FMI, desmonetización del oro; legaliza el no sistema.
\item \textbf{¿Por qué Plaza y luego Louvre?} Plaza para bajar el dólar sobrevalorado; Louvre para frenar su caída.
\item \textbf{¿Qué es Bretton Woods II?} Centro deficitario y periferia con moneda infravalorada que acumula reservas en dólares; estable mientras haya trabajo excedente.
\item \textbf{¿Corrigen los aranceles el déficit?} Poco: $\mathrm{CC}\equiv S-I$; si no cambian el ahorro y la inversión, se aprecia el dólar o se desvía el comercio.
\item \textbf{¿Puede el renminbi sustituir al dólar?} No sin convertibilidad de la cuenta de capital ni mercados profundos; hoy, 2,1\,\% de las reservas.
\item \textbf{¿Qué cambió 2008?} G-20 como foro principal, FSB, Basilea III, \emph{swaps} de la Fed.
\item \textbf{¿Qué reforma del SMI?} Red de seguridad, multipolaridad o DEG supranacional, de más a menos realista.
\item \textbf{¿Qué es el no sistema?} \autor{Williamson} (1976): sin reglas cambiarias comunes ni activo de reserva definido; el orden lo dan los objetivos de inflación y el dólar.
\item \textbf{¿Por qué los emergentes acumulan reservas?} Seguro tras las crisis de los 90 y efecto de mantener la moneda barata (Bretton Woods II); coste: rentabilidad baja y esterilización.
\item \textbf{¿Qué es el privilegio exorbitante?} EE.\,UU. se endeuda en su moneda a tipos bajos e invierte fuera con más rendimiento (\autor{Gourinchas}-\autor{Rey}).
\item \textbf{¿Qué falla hoy en la OMC?} Apelación bloqueada desde 2019 (MPIA como parche), Doha sin cerrar, consenso imposible (CM14), aranceles unilaterales fuera de la NMF.
\item \textbf{¿Por qué la flotación no aísla del todo?} Ciclo financiero global: la Fed mueve precios de activos, crédito y flujos (\autor{Rey}, 2013).
\item \textbf{¿Qué fue el «gran colapso del comercio»?} 2009: comercio $-10$\,\% o más frente a PIB $-0{,}4$\,\%: duraderos, CGV y financiación del comercio.
\end{enumerate}

\begin{trampa}
\begin{itemize}
\item Smithsonian: 38\,\$/onza y bandas $\pm2{,}25$\,\% (antes $\pm1$\,\%); flotación en marzo de 1973, no en 1971.
\item Jamaica en 1976; la Segunda Enmienda entra en vigor en 1978.
\item Plaza (1985): G-5, depreciar; Louvre (1987): G-6 (Italia no firma), estabilizar.
\item México 1994: crisis mixta, no de primera generación pura.
\item Londres 2009: recursos del FMI \emph{triplicados} (750.000 M\$), no duplicados.
\item La Fed subió de 1\,\% a 5,25\,\% entre 2004 y junio de 2006.
\item El Supremo anuló los aranceles IEEPA (feb. 2026), no los de la sección 232.
\item «Pilares» son las instituciones de 1944; «planos», los del SEI.
\end{itemize}
\end{trampa}

\end{columnas}
\end{document}
